% Source: https://artofproblemsolving.com/community/c3h2173912_circular_reasoning
% Stable identity: aops-topic-2173912-q01
\textbf{(a)} The truncation with $n$ digits is
\[0.\underbrace{aa\ldots a}_{n}=a\left(\dfrac{1}{10}+\dfrac{1}{10^2}+\cdots+\dfrac{1}{10^n}\right)=a\cdot\dfrac{1}{10}\cdot\dfrac{1-\dfrac{1}{10^n}}{1-\dfrac{1}{10}}=\dfrac {a}{9}\left(1-\dfrac{1}{10^n}\right),\]
by the formula for a finite geometric series. Its distance from $\dfrac {a}{9}$ is $\dfrac{a}{9\cdot10^n}\le\dfrac{1}{10^n}$, which becomes smaller than any positive number as $n$ grows. So the truncations approach $\dfrac {a}{9}$, and $0.aaa\ldots=\dfrac {a}{9}$.

\textbf{(b)} Reasoning is circular if it uses the statement to be proved. The proof in (a) uses only the definition of an infinite decimal and the finite geometric series, and it works in the same way for every digit $a$; it never assumes anything about $0.999\ldots$. So setting $a=9$ is a legitimate special case, and $0.999\ldots=\dfrac{9}{9}=1$. (The usual shortcut ``$N=0.999\ldots$, $10N-N=9$'' is also valid, but it silently uses that the infinite decimal is a number obeying the usual rules of arithmetic, which is exactly what the limit definition provides.)
